\documentclass[a4paper]{article}
\usepackage{MMVII}

\newcommand{\TheCmd}{ExportSimulMesIm}

\begin{document}

\CmdTitle
%\maketitle

\tableofcontents

%=============================================================================

\SecDescr

This command generates \emph{artificial} correspondences between ground points
and image points for a given set of sensors. It can be used mainly in programs
external to \MMVII for different purposes:

\begin{itemize}
  \item check in a program that an implemented sensor is consistent with a
        given \MMVII sensor;
  \item create a sensor identical to an \MMVII one, using some reverse
        engineering approach;
  \item transform a sensor format, using some bundle adjustment techniques.
\end{itemize}

%=============================================================================

\section{Generated result}

%  -  -  -  -  -  -  -  -  -  -  -  -  -  -  -  -  -  -  -  -  -  -  -  -  -

\subsection{Files and localisation}

A file is generated for each image. This file is located in the folder :

\begin{itemize}
 \item \FileSty{MMVII-PhgrProj/SIMUL/ExportSimulMesIm/}. 
\end{itemize}


It can be generated in any
serialization format supported by \MMVII (\FileSty{xml, json, csv \dots}).
If not specified, it uses the format associated with the profile, but this can
be changed using the \texttt{TypeS} parameter.

%  -  -  -  -  -  -  -  -  -  -  -  -  -  -  -  -  -  -  -  -  -  -  -  -  -

\subsection{Data generated}

The level of detail depends on the type of sensor. In the general case, it
gives only the ground points and their image projections%
\footnote{Bundles may be added later.}.
In the case of central perspective, it gives all the intermediate computations:

\begin{itemize}
  \item \textbf{Ground:} ground point $P_g$ in world coordinates;
  \item \textbf{Cam:} ground point $P_c$ in camera coordinates,
        $P_c = R^t (P_g - C)$;
  \item \textbf{Proj:} projection of $P_c$, which is the result of a mapping
        $\pi : \mathbb{R}^3 \to \mathbb{R}^2$. For a pinhole camera, the
        formula is $\pi(x,y,z) = \frac{(x,y)}{z}$; for a fisheye, there are
        several variants to avoid infinite coordinates;
  \item \textbf{Dist:} application of the distortion;
  \item \textbf{Im:} use of the focal length and the principal point to get
        pixel coordinates.
\end{itemize}

%  -  -  -  -  -  -  -  -  -  -  -  -  -  -  -  -  -  -  -  -  -  -  -  -  -

\subsection{Generated points}

The $3$D points used for creating the correspondences are $100\%$ artificial.
They are generated using the following algorithm:

\begin{itemize}
  \item we generate points $p_{i,j}$ on an $N \times N$ regular grid on the
        sensor;
  \item for each point $p_{i,j}$, we generate $Q$ $3$D points at different
        depths on the bundle issued from $p_{i,j}$.
\end{itemize}

%=============================================================================

\section{Mandatory and optional parameters}

%  -  -  -  -  -  -  -  -  -  -  -  -  -  -  -  -  -  -  -  -  -  -  -  -  -

\subsection{Mandatory parameters}

The two mandatory parameters are:

\begin{itemize}
  \item the set of images to which it applies;
  \item the localisation of the orientations; each image must have an
        orientation, which can be of any type (central perspective,
        pushbroom \dots); the result may vary with the type.
\end{itemize}

%  -  -  -  -  -  -  -  -  -  -  -  -  -  -  -  -  -  -  -  -  -  -  -  -  -

\subsection{Optional parameters}

The optional parameters are:

\begin{itemize}
  \item \textbf{NbPlani, NbDepth}: fix the number of points on the grid and
        the number of depths;
  \item \textbf{TypeS}: fixes the export format (\FileSty{xml, json \dots});
  \item \textbf{OnlyGrIm}: if true, for central perspective, export only
        ground points and images.
  % \item \textbf{EpsilonCheck}: tolerance when we check the image projection;
\end{itemize}

\end{document}
